\documentclass[11pt]{article}
\usepackage[margin=1in]{geometry}
\usepackage{parskip}
\usepackage[hidelinks]{hyperref}
\usepackage{microtype}
\usepackage{amsmath,amssymb}
\begin{document}
\thispagestyle{empty}

\begin{flushright}
College of Science and Engineering\\
James Cook University, Australia\\
\href{mailto:nuonan.ouyang@my.jcu.edu.au}{nuonan.ouyang@my.jcu.edu.au}\\
13 September 2026
\end{flushright}

\vspace{4pt}
Editor-in-Chief and Editorial Team\\
IEEE Transactions on Network and Service Management

\vspace{6pt}
Dear Editor-in-Chief and Editorial Team,

Please consider our new manuscript, ``When Learned IDS Model Selection Collapses to a Static Choice: A Measurement Study at the Edge,'' following the 30 May 2026 decision on prior manuscript TNSM-2026-10776 (``Energy-Aware Intrusion Detection System with Q-learning-Based Dynamic Scheduler on Lightweight IoT Devices''). The decision invited a new manuscript that fully considers the prior reviews; the accompanying response maps each prior comment to the rebuilt evidence and current manuscript.

The submission offers three interlocking, portable contributions to network and service management, together with a physical measurement study that instantiates them at the edge.

\textbf{(1) A pre-deployment reward-audit framework.} We prove that under a linear scalar utility with fixed detection weight, a strict pairwise action reversal on a support region exists if and only if the implied exchange threshold $\tau_{ab}=\Delta C_{ab}/w_D$ lies strictly between the attainable minimum and maximum of the detection-difference $\Delta D_{ab}(\cdot)$ (Theorem~1). Supplementary Algorithm~S1 turns the theorem into a validation-only audit that a deployment engineer can run at policy-freeze time using declared weights, measured cost profiles, and validation-support scores. Under the executed candidate-range normalization, the audit correctly predicts the LightLR collapse of both learned selectors before any test outcome is inspected---this is the paper's core theoretical contribution and is directly applicable to any TNSM study whose learned orchestrator concentrates on a single realized action.

\textbf{(2) An evaluation protocol that separates detector choice from adaptive value.} Its five instruments---paired physical blocks, matched-static-action attribution, prospectively frozen factorial workloads, joint raw and drift-adjusted power endpoints, and full occupancy plus off-support action reporting---together isolate what a learned selector adds \emph{beyond} the schedule it realizes. We consolidate the protocol into a practitioner-facing six-item decision checklist in the Discussion, so the framework is usable by TNSM readers who are not measurement specialists.

\textbf{(3) A positive service-management result via a simple causal rule.} Under a prospectively frozen $50/100/4000\text{ rows/s}\times 0.10/0.50/0.90$ factorial (R2), the causal DeadlineGuard rule cuts energy against Static-MedRF by 43.714~J per 500-second run (95\% paired CI 41.418--46.010~J; 2.98\%) and eliminates all 512/5,000 deadline misses at a 1.373-pp F1 cost. The finding supports a broader design principle: when the switch criterion is expressible from a causal telemetry channel, a simple threshold rule is often preferable to a learned selector on the same channel, both operationally and evaluatively. The finding sits naturally in TNSM's remit because it changes the practical guidance operators receive about learned versus rule-based orchestration under SLO stress.

The manuscript also documents unfavorable outcomes for the learned policies: both realize a static LightLR schedule on the observed support, and when the R2 factorial pushes Tabular-Q off the training support, its zero-initialized ties promote a TinyDT default with 67.76\% FPR. These outcomes are retained because they anchor the safety principle that discrete controllers must audit their default action assignments, not only their occupancy on visited states.

Relative to the prior submission, we rebuilt the data lineage, detector evaluation, scheduler implementation, and physical evidence, and added DQN as the stronger DRL comparator, explicit causal state/reward definitions, reward and state sensitivity, learning curves, matched physical statistics, isolated scheduler-overhead measurements, a validation-only TinyDT threshold diagnostic, and an expanded related-work discussion covering conditional cascades and inference-service management. Multi-node scalability remains explicitly outside the evidence boundary and is not claimed.

All known failed or invalid attempts and authorized replacements remain archived, executed-source identities are hash-bound, and the evidence packages pass CRC, manifest, source-binding, and test-leakage checks. The reproducibility archive is intended for release under a DOI upon acceptance. We confirm that the submission is original, is not under consideration elsewhere, and is submitted with the approval of all authors.

Thank you for your consideration.

\vspace{18pt}
Sincerely,

\vspace{34pt}
Nuonan Ouyang\\
on behalf of all authors

\end{document}
